\documentclass{plantillaPPT}

\titulo{12}{Fin de Polares e Inicio de Series y Sucesiones}

\begin{document}
\primeraDiapo

\begin{frame}{Polares}
\framesubtitle{Práctica 12}

1. Considere las siguientes curvas en coordenadas polares:
\[r_1 = 3\cos(\theta) \qquad r_2=2+\cos(\theta)\]\\[5ex]
\((a)\) Determinar, si existen, los puntos de intersección entre las curvas.\\[3ex]
\((b)\) Mostrar que ambas curvas son simétricas con respecto al eje polar.

\end{frame}

\begin{frame}{Polares}
\framesubtitle{Práctica 12}
\((c)\) Calcular el área de la región exterior a \(r_1\) e interior a \(r_2\).\\[2ex]

\begin{center}
    \includegraphics[width=0.5\textwidth]{C2 2025-1/imagenes/ayud12img.png}
\end{center}
    
\end{frame}

\begin{frame}{Verdadero o Falso}
\framesubtitle{Práctica 12}

2. Decidir si las siguientes afirmaciones son verdaderas o falsas. Justificar su respuesta.\\[3ex]
\((a)\) Sea \((a_n)_{n\in\mathbb{N}}\) la sucesión definida por: \[a_1=1 \text{ y } a_{n+1}=\sqrt{3+2a_n}\] 
Para \(n\geq1\). \\[2ex]
Si \((a_n)_{n\in\mathbb{N}}\) es monótona creciente y, \(\forall n>1, n\in\mathbb{N}\), se cumple que \(2<a_n<3\), entonces \(\lim_{n\to+\infty}a_n=3\).

\end{frame}

\begin{frame}{Verdadero o Falso}
\framesubtitle{Práctica 12}

\((b)\) Si la sucesión \((a_x)_{n\in\mathbb{N}}\) es convergente, entonces la serie \\ \(\displaystyle\sum_{n=1}^{\infty}a_n\) es convergente.\\[4ex]
\((c)\) Si la sucesión \((a_n)_{n\in\mathbb{N}}\) es divergente, entonces la serie \\\(\displaystyle\sum_{n=1}^\infty \frac{1}{a_n}\) es divergente.
    
\end{frame}

\begin{frame}{Cálculo de Series}
\framesubtitle{Práctica 12}

3. Calcule el valor de las siguientes sumas:\\[5ex]

\((a) \displaystyle\sum_{n=1}^{\infty}\frac{2}{16n^2-8n-3}\) \\[3ex]

\((b) \displaystyle\sum_{n=1}^{\infty}\frac{2}{4n^2-1}\)

\end{frame}

\end{document}